\section{Recomanacions}

\subsection{Interstellar}

[draft] peli interestellar

\subsection{Ant-Man}

Ant-Man juega con ideas reales de la física cuántica, pero las lleva bastante más allá de donde la ciencia permite llegar. Partículas Pym, cambios de tamaño y un “reino cuántico” sirven como excusa para montar la aventura. No es una clase de física, desde luego, pero sí una forma divertida de acercarse a conceptos científicos y descubrir que, a veces, Hollywood interpreta la física con bastante imaginación.

\subsection{Viatge al·lucinant, d'Isaac Asimov}

[draft] libro viaje alucinante

Viaje alucinante, de Isaac Asimov, imagina una tripulación miniaturizada viajando por el interior del cuerpo humano para operar un coágulo cerebral. La miniaturización es pura ciencia ficción y físicamente hace aguas por bastantes sitios, pero la anatomía y muchos de los problemas que encuentran —circulación, pulmones, sistema inmunitario— sí parten de ideas reales.

Como ciencia no conviene tomárselo demasiado al pie de la letra. Como puerta de entrada a la biología y a preguntarse “¿qué pasaría si pudiéramos hacer esto?”, funciona estupendamente.
